% ---- pg-417 （书页 405） ----

\begin{mvtable}
\phantomsection\label{tab:9.10}\mvtabcap{表 9.10\quad 鸡的骨骼数据的因子分析}
\mvsub{主成分法}
\small
\begin{tabular}{lccccccc}
\toprule
& \multicolumn{3}{c}{因子载荷估计} & \multicolumn{3}{c}{旋转因子载荷估计} & \\
\cmidrule(lr){2-4}\cmidrule(lr){5-7}
变\quad 量 & $F_1$ & $F_2$ & $F_3$ & $F_1^{*}$ & $F_2^{*}$ & $F_3^{*}$ & $\tilde{\psi}_i$\\
\midrule
1. 颅骨长度 & 0.741 & 0.350 & 0.573 & 0.355 & 0.244 & 0.902 & 0.00\\
2. 颅骨幅宽 & 0.604 & 0.720 & $-0.340$ & 0.235 & 0.949 & 0.211 & 0.00\\
3. 股骨长 & 0.929 & $-0.233$ & $-0.075$ & 0.921 & 0.164 & 0.218 & 0.08\\
4. 胫骨长 & 0.943 & $-0.175$ & $-0.067$ & 0.904 & 0.212 & 0.252 & 0.08\\
5. 肱骨长 & 0.948 & $-0.143$ & $-0.045$ & 0.888 & 0.228 & 0.283 & 0.08\\
6. 尺骨长 & 0.945 & $-0.189$ & $-0.047$ & 0.908 & 0.192 & 0.264 & 0.07\\
\midrule
\makecell[l]{所解释的(标准化)\\样本总方差的\\累积比例}
 & 0.743 & 0.873 & 0.950 & 0.576 & 0.763 & 0.950 & \\
\bottomrule
\end{tabular}

\vspace{10pt}
\mvsub{极大似然法}
\begin{tabular}{lccccccc}
\toprule
& \multicolumn{3}{c}{因子载荷估计} & \multicolumn{3}{c}{旋转因子载荷估计} & \\
\cmidrule(lr){2-4}\cmidrule(lr){5-7}
变\quad 量 & $F_1$ & $F_2$ & $F_3$ & $F_1^{*}$ & $F_2^{*}$ & $F_3^{*}$ & $\hat{\psi}_i$\\
\midrule
1. 颅骨长度 & 0.602 & 0.214 & 0.286 & 0.467 & 0.506 & 0.128 & 0.51\\
2. 颅骨幅宽 & 0.467 & 0.177 & 0.652 & 0.211 & 0.792 & 0.050 & 0.33\\
3. 股骨长 & 0.926 & 0.145 & $-0.057$ & 0.890 & 0.289 & 0.084 & 0.12\\
4. 胫骨长 & 1.000 & 0.000 & $-0.000$ & 0.936 & 0.345 & $-0.073$ & 0.00\\
5. 肱骨长 & 0.874 & 0.463 & $-0.012$ & 0.831 & 0.362 & 0.396 & 0.02\\
6. 尺骨长 & 0.894 & 0.336 & $-0.039$ & 0.857 & 0.325 & 0.272 & 0.09\\
\midrule
\makecell[l]{所解释的(标准化)\\样本总方差的\\累积比例}
 & 0.667 & 0.738 & 0.823 & 0.559 & 0.779 & 0.823 & \\
\bottomrule
\end{tabular}
\end{mvtable}

旋转极大似然因子载荷与用主成分法生成的因子载荷对第一因子是一致的, 但对第二和第三因子不一致. 对极大似然法, 第二因子是代表头部大小. 第三因子含义不清, 而且可能是不需要的.

对保留三个或更少因子的进一步支持, 由从极大似然估计得到的残差阵提供:
\[
\mathbf{R}-\hat{\mathbf{L}}_z\hat{\mathbf{L}}_z'-\hat{\boldsymbol{\Psi}}_z=
\left[\begin{array}{rrrrrr}
0.000& & & & & \\
-0.000&0.000& & & & \\
-0.003&0.001&0.000& & & \\
0.000&0.000&0.000&0.000& & \\
-0.001&0.000&0.000&0.000&0.000& \\
0.004&-0.001&-0.001&0.000&-0.000&0.000
\end{array}\right]
\]

该矩阵的全部元都很小. 我们将继续探索此例的 $m=3$ 因子模型. $m=2$ 因子模型在练习 9.10 中考虑.

从带有极大似然估计旋转的\hyperref[eq:9-58]{式(9-58)}产生的对因子 1 和因子 2 的因子得分, 示于\hyperref[fig:9.5]{图 9.5} 中. 该图可让我们确认那些出于某种原因与其余观测值不一致的观测值. 可能的离群值在图上已圈出.
